% \lettrine{}{} command typesets a drop cap in the document

% second brace set formats text placed inside as small caps
% (same thing as doing \lettrine{L}\textsc{orem ipsum ...}

\chapter{Precio (Redención?)}
\lettrine{L}{orem ipsum dolor sit amet}. []Post-final]

Parece que ya todo pasó. Los tres amigos [nuestro prota, una muchacha de la misma edad y un muchacho más chico, casi un niño] se despiden y se van cada quien para su casa.

Al día siguiente, por la mañana, nuestro protagonista se levanta y empieza a tomar una ducha. Sin embargo, no puede dejar de pensar en los acontecimientos del día anterior. El [ente] parecía vencido [aplacado], pero de algún modo las cosas habían resultado demasiado fáciles. Y de algún modo, la actitud de [la muchacha] le había parecido bastante extraña, incluso al final del todo. No parecía disfrutar de forma sincera el triunfo y parecía más bien distante y pensativa.

De pronto, le llegó fulminantemente la revelación. Veladamente, con aquellas últimas palabras que dijo, ella había pactado con el [ente] y se había ofrecido para ir sola al día siguiente. Consciente de pronto del peligro que corría su amiga, y de que probablemente había salido desde temprano en la madrugada para llegar al [lugar sagrado], su cara se transforma en terror puro. [Nuestro prota] se medio enjuaga y se medio seca y sale a la habitación a vestirse. -- Ella no sabe de qué cosas [el ente] es capaz --, piensa. Al menos, no tanto como él [quien fue el primero a quien la entidad empezó a molestar]. Lamenta no haberle dicho todo desde el principio mientras se mete lo más rápido que puede en un overol de mezclilla. No hay tiempo para ponerse las botas. Se arremanga el oberol para no tropezarse y corre descalzo hasta donde está el caballo, se monta y cabalga cerro arriba a toda velocidad. No sabe qué va a hacer al llegar, pero no va a dejar que ella se enfrente sola a eso, ¡No! Una sola persona no tendría ninguna oportunidad. Sería prácticamente una ofrenda. ¿Acaso ella piensa scrificarse para de algún modo protegerlos a ellos?

Llegado a la cerca desmonta y corre colina abajo. El primer trecho cubierto sólo de pasto es suave, pero llegado a la sombra de los árboles, donde el suelo está cubierto se agujas de [pino], le lastima pisar las pequeñas piñas [no sé de qué árbol son ni si tienen nombre en especial. Investigar]. No está acostumbrado a andar descalzo, pero con el dolor y todo no disminuye su velocidad. Al fin, al darse cuenta de que ya está muy cerca de las grandes rocas, se para casi en seco y empieza a avanzar cautelosamente, pendiente de cualquier ruido que pueda escuchar, pero al mismo tiempo sin poder controlar todavía la agitación de su respiración y de su corazón. Le angustia pensar qué se va a encontrar. El terror de volver a enfrentarse al [ente] no palidece ante el de haber llegado demasiado tarde.

Al asomarse al centro del conjunto de piedras, ve a [la muchacha] sentada sobre una de ellas, que tiene casi la altura y la forma de una pequeña banca. Tiene un vestido blanco de flores, el pelo recogido en una trenza, y la mirada hacia abajo, pensativa. El cielo está claro, pero aún es demasiado temprano para que los árboles proyecten alguna sombra. En el fresco del amanecer se respira paz, que en contraste con la agitación del muchacho se percibe casi irreal. En el silencio de la mañana en el monte, ella lo escucha acercarse aún desde lejos y levanta la mirada. Se ve totalmente sorprendida, pero nada asustada.

-- Hola -- dice ella -- ¿Qué haces aquí?. -- Lo mismo digo -- dice él -- ¿Dónde está [el ente]?. --Se fué-- dice ella. [Explica por qué fue ella sola, que estaba preocupada por él, a quien ama, y por el amigo más chico, que ven como, o es, un hermano pequeño, y no quería perderlos. Cómo el ente apareció, pero ella lo aplacó, parece que finalmente y para siempre, gracias al poder que ella tenía, heredado de su familia. Para esto, confiesa la relación entre la entidad y su familia, que antes no estaba clara.] [Nuestro prota] escucha atentamente el relato. Su angustia, que se mantenía apretada bajo su pecho debido a la urgencia, empieza a disiparase, y a propagarse desde ahí. Llega hasta sus brazos y piernas, haciéndole sentirlos guangos [buscar una mejor palabra] y sube hasta que se le atora en el cogote, haciéndole un nudo en la garganta. Sabe que no puede hablar sin que se le quiebre la voz, pero aún así dice: -- Pensé que no te iba a volver a ver --.

Ella baja de nuevo la mirada. --Perdóname-- dice gravemente y sin levantar la mirada. -- Ven, siéntate--. Se hace a un lado para hacerle espacio de sentarse sobre la misma roca. Se quedan en silencio durante un rato, luego ella le toma la mano. Él la voltea a ver a la cara. Sus rostros están demasiado cerca, y ella puede sentir el calor de la agitación en el de él. De pronto, ella le pone la mano sobre la mejilla y se besan. [Él se estremece y siente que la angustia 'is flushed away'. Ella también parece quitarse un enorme peso]
-- ¿Entonces te viniste corriendo desde el pueblo?
-- Me vine cabalgando desde mi casa, sólo corrí desde la cerca.
Ella le sacude el cabello y le dice:
-- Tienes el pelo lleno de basuras de árbol, llevas un tirante desabrochado y tienes sangre en tu pie. ¿Por qué te viniste sin zapatos y sin camisa?
De pronto, él se siente ligeramente avergonzado por mostrarse en un estado tan lamentable, al mismo tiempo que se da cuenta de que tiene los brazos y la cara llenos de arañazos. Se ruboriza y empieza a justificarse.
-- Pues es que no sabía que ibas a regresar aquí. Caí en la cuenta mientras me estaba bañando. Salí corriendo de mi casa y sólo me puse este overol. De seguro me abroché mal el tirante...
Ella sonríe nerviosa.
-- ¿Entonces te viniste a "ráiz"?
-- A "ráiz" --. Confirma él sonriendo, y se levanta ligeramente la tela del overol al lado de los botones laterales, que vienen desabrochados, dejando al descubierto un pedacito de piel a la altura de la cadera. Ambos ríen. Luego hay unos segundos de silencio. Ella le pone su mano sobre el botón del tirante que sí está abrochado, acerca su cara a la de él, y le dice en voz baja: --Quítate el otro--. El botón se suelta, dejando el peto caer, al tiempo que los dos comienzan a besarse nuevamente. Ella lo empuja hacia atrás y él se deja caer.

Lo que sigue no lo puedo contar. Baste saber que ese día los dos entraron cabalgando juntos al pueblo, donde ya los estaban esperando preocupados. Las familias esaban escandalizadas de que hubieran pasado todo el día juntos, ¡y solos! y quisieron presionarlos para casarse, pero no tuvieron que presionar demasiado. Se casaron en la fiesta patronal tres meses después, y a los seis meses nació su primer hijo. [Una buena observación para decir que "vivieron felices por siempre" y tal vez conectar con la época actual, también, decir qué pasó con el tercer amigo del grupo].